\begin{figure}[t]
\centering
\resizebox{0.78\textwidth}{!}{
\begin{tikzpicture}[
box/.style={
rectangle,draw,rounded corners,
align=center,
minimum width=5.4cm,
minimum height=1.15cm
},
arrow/.style={->,thick}
]

\node[box] (harm) {Harmonic geometry\\equilibrium structure + inertia tensor};
\node[box,below=1.3cm of harm] (muone) {$\mu_1=\partial I^{-1}/\partial Q$\\first geometric rotation--vibration coupling};
\node[box,below=1.3cm of muone] (linear) {Linear tensor sector\\standard quartic + standard sextic};
\node[box,below=1.3cm of linear] (cubic) {Cubic force constants\\explicit anharmonic input};
\node[box,below=1.3cm of cubic] (nonlinear) {First nonlinear layer\\first breakdown of linear transport};

\draw[arrow] (harm.south) -- (muone.north);
\draw[arrow] (muone.south) -- (linear.north);
\draw[arrow] (linear.south) -- (cubic.north);
\draw[arrow] (cubic.south) -- (nonlinear.north);

\end{tikzpicture}
}

\caption{
Conceptual hierarchy revealed by the tensor formulation. Harmonic
geometry defines the inertia tensor, its first derivative $\mu_1$
generates the minimal linear tensor sector governing standard quartic
and standard sextic centrifugal distortion constants, and cubic force
constants introduce the first genuinely nonlinear layer.
}

\label{fig:ceditt3_hierarchy}
\end{figure}
